\section{Well-foundedness is Essential}
\label{sec:lob-essential}

The abstract framework of Section~\ref{sec:abstract} carries one
distinctive axiom: \texttt{fs\_lt\_wf}, the well-foundedness of the
strict refinement relation.  Every other axiom of
\texttt{ForcingStructure} is an order-theoretic structural one:
reflexivity, transitivity, compatibility.  Well-foundedness is the
exception, and the only axiom that uses any form of induction
principle to state.  It is natural to ask whether it is essential:
could we drop it and still get a sensible Löb's rule for some other
reason?

The answer is no, and we now make this precise.  We will exhibit a
``would-be'' iProp framework over the integers $\mathbb{Z}$, in which
the carrier and the order satisfy every part of
\texttt{ForcingStructure} \emph{except} well-foundedness.  All the
constructions of Sections~\ref{sec:prop} and~\ref{sec:abstract} go
through: we can build \texttt{iPropZ}, \texttt{iImplZ}, \texttt{iLaterZ}.
But Löb's rule fails.

\paragraph{The setup.}
We work with the integers under $\le$ and $<$ in the obvious way.
Define \texttt{iPropZ} as a record of downward-closed predicates on
$\mathbb{Z}$:

\begin{lstlisting}
Record iPropZ : Type := MkPropZ {
  iholdZ : Z -> Prop;
  imonoZ : forall p q, (q <= p)%Z -> iholdZ p -> iholdZ q
}.

Definition ientailsZ (φ ψ : iPropZ) : Prop :=
  forall p, iholdZ φ p -> iholdZ ψ p.
\end{lstlisting}

\noindent
The connectives follow the same definitions as before.  We give the
ones we need explicitly:

\begin{lstlisting}
Definition iFalseZ : iPropZ :=
  MkPropZ (fun _ => False) (fun _ _ _ H => H).

Definition iImplZ (φ ψ : iPropZ) : iPropZ :=
  MkPropZ
    (fun p => forall q, (q <= p)%Z ->
                          iholdZ φ q -> iholdZ ψ q)
    (iImplZ_mono φ ψ).

Definition iLaterZ (φ : iPropZ) : iPropZ :=
  MkPropZ
    (fun p => forall p', (p' < p)%Z -> iholdZ φ p')
    (iLaterZ_mono φ).
\end{lstlisting}

\noindent
The definitions are the same as in Section~\ref{sec:abstract}, with
$\mathtt{fs\_cond} := \mathbb{Z}$, $\mathtt{fs\_le} := \le_\mathbb{Z}$,
and $\mathtt{fs\_lt} := <_\mathbb{Z}$.  The order axioms hold: $\le$
is a preorder, $<$ is contained in $\le$, and $<$ is compatible with
$\le$ on the right.  Only well-foundedness fails: there is no
minimum and so no base case for an inductive argument.

\paragraph{The counterexample.}
Consider $\triangleright \bot$, the iProp asserting that
\texttt{False} holds at every strict predecessor.  We claim this is
forced at \emph{no} condition.  Indeed, $p \Vdash \triangleright \bot$
would say ``for every $p' < p$, $\bot$ holds at $p'$'', i.e., ``for
every integer strictly less than $p$, \texttt{False} is true''.
There is at least one such integer, namely $p - 1$, and
\texttt{False} fails there.  So $\triangleright \bot$ is the empty
iProp:

\begin{lstlisting}
Lemma lob_premise_at p :
  p ⊩ (iImplZ (iLaterZ iFalseZ) iFalseZ).
Proof.
  intros q Hqp Hlater.
  apply (Hlater (q - 1)); lia.
Qed.
\end{lstlisting}

\noindent
The lemma actually shows something more: at every $p$, the
implication $\triangleright \bot \to \bot$ is forced.  The reason is
the same as just above.  The antecedent $\triangleright \bot$ is
empty, so any implication from it is vacuously forced.  Spelled out:
at $p$, $p \Vdash (\triangleright \bot \to \bot)$ unfolds to ``for
every $q \le p$, if $q \Vdash \triangleright \bot$ then $q \Vdash
\bot$''. For every $q$, the antecedent $q \Vdash \triangleright \bot$
is itself false (apply \texttt{Hlater} at $q - 1$, contradiction), so
the conditional is vacuously satisfied.

The conclusion of Löb's rule says $\bot$ is forced at every
condition.  This obviously fails:

\begin{lstlisting}
Lemma lob_conclusion_fails :
  ~ (forall p, iholdZ iFalseZ p).
Proof. intros H. exact (H 0). Qed.
\end{lstlisting}

\noindent
Combining the two: the entailment $(\triangleright \bot \to \bot)
\vDash \bot$ is the statement ``for every condition, if
$(\triangleright \bot \to \bot)$ is forced there then $\bot$ is forced
there''. The premise is forced everywhere, the conclusion is forced
nowhere, so the entailment fails.

\begin{lstlisting}
Theorem lob_fails_over_Z :
  ~ (iImplZ (iLaterZ iFalseZ) iFalseZ ⊨⊢ iFalseZ).
Proof.
  intros Hent. apply lob_conclusion_fails.
  intros p. apply Hent. apply lob_premise_at.
Qed.
\end{lstlisting}

This concludes the counterexample: there is a forcing-like structure
over $\mathbb{Z}$ in which every other piece of the abstract
framework of Section~\ref{sec:abstract} goes through, but Löb's rule
is not derivable.

\paragraph{From example to characterization.}
The counterexample shows that well-foundedness is necessary on at
least one structure.  We now upgrade ``at least one'' to ``every'',
mechanizing a precise iff.  Fix a ``pre-structure'': a preorder
$(\mathtt{cond}, \le)$ with a strict relation $\mathtt{lt}$ satisfying
all the order-theoretic axioms of \texttt{ForcingStructure} except
well-foundedness.  Build $\mathtt{iPropPre}$,
$\mathtt{iImplPre}$, $\mathtt{iLaterPre}$ exactly as in
Section~\ref{sec:abstract}.  Say that \emph{L\"ob holds} for the
pre-structure if, for every $\mathtt{iPropPre}$ $\varphi$, the
entailment
$(\triangleright \varphi \to \varphi) \vDash \varphi$ is derivable.
Then:

\begin{lstlisting}
Theorem lob_iff_wf : lob_holds <-> well_founded lt.
\end{lstlisting}

\noindent
The forward direction $(\Leftarrow)$ is the abstract \texttt{iLob}
proof of Section~\ref{sec:abstract}, replayed in the pre-structure.
The backward direction $(\Rightarrow)$ is genuinely new.  Define the
\emph{forcing-accessibility} predicate
$\varphi_{\mathrm{acc}}\, p := \mathtt{Acc}\, \mathtt{lt}\, p$. This is
a valid $\mathtt{iPropPre}$, because $\mathtt{Acc}$ is downward-closed
along $\le$.  Its L\"ob premise
$(\triangleright \varphi_{\mathrm{acc}} \to \varphi_{\mathrm{acc}})$ is
forced at \emph{every} condition: $\mathtt{Acc}$'s introduction rule
says exactly ``every strict predecessor accessible $\Rightarrow$ this
condition accessible''.  The assumed universal L\"ob then concludes
$\varphi_{\mathrm{acc}}$ holds everywhere, i.e., every condition is
accessible, i.e., $\mathtt{lt}$ is well-founded.

The statement is the step-indexed Hoare logic form of a classical
soundness--completeness pairing for the G\"odel--L\"ob axiom over
transitive Kripke frames (recalled below).  What is novel is the
combination: stated as a property of an abstract step-indexed forcing
framework, the witness is a single iProp internal to that framework
(forcing-accessibility), and the result is mechanized in Coq.  To our
knowledge no equivalent characterisation has been recorded for any
extant step-indexed Hoare logic.

\paragraph{Connection to G\"odel--L\"ob.}
The result is the step-indexed analog of a classical observation in
modal logic.  The G\"odel--L\"ob provability logic \textsf{GL} is
sound and complete for the class of transitive Kripke frames whose
accessibility relation is converse well-founded
\cite{smorynski1985,boolos1993}.  Frames where the accessibility
relation is \emph{not} converse well-founded falsify the Löb axiom.
The integers under $<$ are the canonical example.  Our
counterexample is the obvious instantiation of that classical fact
into the step-indexed Hoare logic setting: forcing conditions are
the worlds, strict refinement is the converse of the accessibility
relation, $\triangleright$ is the necessity modality, and Löb's rule
is the Löb axiom.  Well-foundedness of the strict refinement is
exactly converse well-foundedness of the accessibility relation.

\paragraph{Methodological upshot.}
The ``ramification'' in ramified forcing is not a presentation
convenience to be eliminated.  Shoenfield's unramification of
set-theoretic forcing showed that for forcing the explicit ordinal
hierarchy can be absorbed into a single recursion on rank, because
the relevant induction principle is already available at the
meta level.  In the step-indexed setting, the analogous move is
\emph{not} available, because Löb's rule depends on a property of
the forcing notion (well-foundedness) that cannot be supplied by any
amount of recursion at the meta level: it has to be a property of
the conditions themselves.  The well-foundedness is constitutive of
the modal logic and not eliminable.

\paragraph{Practical upshot.}
The choice of step-index domain decides whether L\"ob's rule is
available.  Building a step-indexed semantics on a non-well-founded
index structure, such as signed counters, bidirectional approximation
indices, or anything descending without bound, yields every iProp
connective and a derivable $\triangleright$-monotonicity rule, but
L\"ob's rule fails: $(\triangleright \bot) \to \bot$ becomes
vacuously valid while $\bot$ stays unforced, and every soundness
proof whose predicate is recursive in the program
(Section~\ref{sec:imp-vs-lam}) silently relies on a rule that does
not hold.  The standard candidates all pass the check: $\omega$,
products like $\omega \times \omega$, ordinals up to a bound, and
Transfinite Iris~\cite{spies2021transfinite}.  Our
characterization rules out only the temptation to leave the
well-foundedness implicit.

